\begin{pastexam}[label={px:gk3-110-ds-01}]{高考三級資訊處理 110 年 第一題}
（申論；25 分）

A 為（8×4）矩陣、B 為（4×10）矩陣、C 為（10×3）矩陣、D 為（3×20）矩陣、E 為（20×4）矩陣，
㈠請列出此 5 個矩陣相乘 A×B×C×D×E 所有可能的乘法順序（請用括號表示乘法順序）。（5 分）
㈡請使用 Dynamic Programming（動態規劃）的技巧計算出此五個矩陣相乘 A×B×C×D×E 的最佳乘法順序
（請用括號表示乘法順序），使得五個矩陣相乘所需要花費的乘法數量最少。（15 分）
㈢請列出此五個矩陣相乘所需要花費的最少乘法數量。（5 分）（注意：未說明 Dynamic Programming 的
計算過程，不予計分。）
\end{pastexam}

\begin{pastsoln}
答案：最佳順序 $(A((BC)(DE)))$，最少 $536$ 次乘法。

記號照資結書第~18.2~節：$A_1, \ldots, A_5$ 依序是 A、B、C、D、E，$A_i$ 是 $p_{i-1} \times p_i$，
\[
\langle p_0, p_1, p_2, p_3, p_4, p_5 \rangle = \langle 8, 4, 10, 3, 20, 4 \rangle .
\]
$p \times q$ 乘 $q \times r$ 要 $pqr$ 次純量乘法。

\textbf{㈠ 所有順序}。 $n$ 個矩陣加滿括號的方法有卡特蘭數 $C_{n-1}$ 種，$5$ 個矩陣是
$C_4 = 14$ 種。 依最後一次乘法把鏈切在哪裡分類（切在第 $k$ 個之後，左邊 $k$ 個有 $C_{k-1}$
種、右邊 $5 - k$ 個有 $C_{4-k}$ 種）：$C_0 C_3 + C_1 C_2 + C_2 C_1 + C_3 C_0 = 5 + 2 + 2 + 5 = 14$。
下表列出全部 $14$ 種和各自的乘法次數（次數是㈡的驗證用）：
\begin{center}
\small
\begin{tabular}{llr@{\qquad}llr}
\toprule
最後切點 & 順序 & 次數 & 最後切點 & 順序 & 次數 \\
\midrule
A $|$ BCDE & $A(B(C(DE)))$ & 648 & ABC $|$ DE & $(A(BC))(DE)$ & 552 \\
 & $A(B((CD)E))$ & 1688 & & $((AB)C)(DE)$ & 896 \\
 & $A((BC)(DE))$ & \textbf{536} & ABCD $|$ E & $(A(B(CD)))E$ & 2680 \\
 & $A((B(CD))E)$ & 1848 & & $(A((BC)D))E$ & 1640 \\
 & $A(((BC)D)E)$ & 808 & & $((AB)(CD))E$ & 3160 \\
AB $|$ CDE & $(AB)(C(DE))$ & 1000 & & $((A(BC))D)E$ & 1336 \\
 & $(AB)((CD)E)$ & 2040 & & $(((AB)C)D)E$ & 1680 \\
\bottomrule
\end{tabular}
\end{center}
（最外層的括號省略。）

\textbf{㈡ 動態規劃}。 令 $m[i, j]$ 為算出 $A_i A_{i+1} \cdots A_j$ 最少的乘法次數。 最後一次乘法把
這段切成 $A_i \cdots A_k$ 和 $A_{k+1} \cdots A_j$，兩段各自也要用最省的方法算（最佳子結構），
所以（資結書第~18.2~節）
\[
m[i, j] = \begin{cases}
0, & i = j, \\
\displaystyle\min_{i \le k < j} \bigl\{ m[i, k] + m[k+1, j] + p_{i-1} p_k p_j \bigr\}, & i < j,
\end{cases}
\]
並記下取到最小值的 $k$ 為 $s[i, j]$。 依鏈長 $2, 3, 4, 5$ 的順序填表：

長度 $2$（只有一種切法）：
\begin{itemize}
\item $m[1,2] = 8 \cdot 4 \cdot 10 = 320$；$m[2,3] = 4 \cdot 10 \cdot 3 = 120$；
$m[3,4] = 10 \cdot 3 \cdot 20 = 600$；$m[4,5] = 3 \cdot 20 \cdot 4 = 240$。
\end{itemize}
長度 $3$：
\begin{itemize}
\item $m[1,3] = \min\{\, m[2,3] + 8 \cdot 4 \cdot 3,\; m[1,2] + 8 \cdot 10 \cdot 3 \,\}
= \min\{120 + 96, 320 + 240\} = \min\{216, 560\} = 216$，$s = 1$。
\item $m[2,4] = \min\{\, m[3,4] + 4 \cdot 10 \cdot 20,\; m[2,3] + 4 \cdot 3 \cdot 20 \,\}
= \min\{600 + 800, 120 + 240\} = \min\{1400, 360\} = 360$，$s = 3$。
\item $m[3,5] = \min\{\, m[4,5] + 10 \cdot 3 \cdot 4,\; m[3,4] + 10 \cdot 20 \cdot 4 \,\}
= \min\{240 + 120, 600 + 800\} = \min\{360, 1400\} = 360$，$s = 3$。
\end{itemize}
長度 $4$：
\begin{itemize}
\item $m[1,4]$：$k = 1$：$0 + 360 + 8 \cdot 4 \cdot 20 = 1000$；$k = 2$：$320 + 600 + 8 \cdot 10 \cdot 20 = 2520$；
$k = 3$：$216 + 0 + 8 \cdot 3 \cdot 20 = 696$。 $m[1,4] = 696$，$s = 3$。
\item $m[2,5]$：$k = 2$：$0 + 360 + 4 \cdot 10 \cdot 4 = 520$；$k = 3$：$120 + 240 + 4 \cdot 3 \cdot 4 = 408$；
$k = 4$：$360 + 0 + 4 \cdot 20 \cdot 4 = 680$。 $m[2,5] = 408$，$s = 3$。
\end{itemize}
長度 $5$：
\begin{itemize}
\item $m[1,5]$：$k = 1$：$0 + 408 + 8 \cdot 4 \cdot 4 = 536$；$k = 2$：$320 + 360 + 8 \cdot 10 \cdot 4 = 1000$；
$k = 3$：$216 + 240 + 8 \cdot 3 \cdot 4 = 552$；$k = 4$：$696 + 0 + 8 \cdot 20 \cdot 4 = 1336$。
$m[1,5] = 536$，$s = 1$。
\end{itemize}
整理成表（左：$m[i,j]$；右：$s[i,j]$）：
\begin{center}
\small
\begin{tabular}{c|ccccc}
$m$ & $j=1$ & 2 & 3 & 4 & 5 \\
\hline
$i=1$ & 0 & 320 & 216 & 696 & 536 \\
2 & & 0 & 120 & 360 & 408 \\
3 & & & 0 & 600 & 360 \\
4 & & & & 0 & 240 \\
5 & & & & & 0 \\
\end{tabular}
\qquad
\begin{tabular}{c|cccc}
$s$ & $j=2$ & 3 & 4 & 5 \\
\hline
$i=1$ & 1 & 1 & 3 & 1 \\
2 & & 2 & 3 & 3 \\
3 & & & 3 & 3 \\
4 & & & & 4 \\
\end{tabular}
\end{center}
由 $s$ 表還原順序：$s[1,5] = 1$，切成 $A_1 \mid A_2 \cdots A_5$；$s[2,5] = 3$，切成
$A_2 A_3 \mid A_4 A_5$。 所以最佳順序是
\[
(A((BC)(DE))) .
\]
也就是先算 $BC$（$4 \times 3$）和 $DE$（$3 \times 4$），再相乘得 $4 \times 4$，最後用 A 去乘。

\textbf{㈢ 最少次數}：$m[1,5] = 536$。 照順序加一遍驗證：$BC$ 要 $4 \cdot 10 \cdot 3 = 120$，$DE$ 要
$3 \cdot 20 \cdot 4 = 240$，$(BC)(DE)$ 要 $4 \cdot 3 \cdot 4 = 48$，$A$ 乘上去要 $8 \cdot 4 \cdot 4 = 128$，
合計 $120 + 240 + 48 + 128 = 536$，和㈠的表裡最小的一個相同。 次佳的 $((A(BC))(DE))$ 要 $552$ 次，
由左往右依序乘要 $1680$ 次。
\end{pastsoln}
